\begin{theorem}[window-\texorpdfstring{$6$}{6} 的 \texorpdfstring{$C_3$}{C3} 根字典满足二次矩各向同性与二次能量等分]\label{thm:xi-window6-c3-quadratic-energy-equipartition}
\leanverified{paper\_xi\_window6\_c3\_quadratic\_energy\_equipartition}
沿用定理 \ref{thm:xi-window6-center-exact-splitting-bdry-cyclic} 的 window-\(6\) 记号，并令
\[
\iota_C:X_6^{\mathrm{cyc}}\to \mathbb R^3
\]
为表 \ref{tab:fold6_b3c3_root_dictionary_c3} 给出的 \(C_3\) 根字典。记
\[
L:=\{w\in X_6^{\mathrm{cyc}}:\ \mathrm{wt}(w)=1\},
\qquad
S:=X_6^{\mathrm{cyc}}\setminus L.
\]
则 \(L\) 恰对应六个长根
\[
\{\pm(2,0,0),\ \pm(0,2,0),\ \pm(0,0,2)\},
\]
而 \(S\) 恰对应十二个短根
\[
\{(\pm1,\pm1,0),\ (\pm1,0,\pm1),\ (0,\pm1,\pm1)\}.
\]
并且成立精确二次框架恒等式
\[
\boxed{\;
\sum_{w\in L}\iota_C(w)\iota_C(w)^{\mathsf T}=8I_3,\qquad
\sum_{w\in S}\iota_C(w)\iota_C(w)^{\mathsf T}=8I_3.\;}
\]
从而
\[
\boxed{\;
\sum_{w\in X_6^{\mathrm{cyc}}}\iota_C(w)\iota_C(w)^{\mathsf T}=16I_3.\;}
\]
等价地，
\[
\boxed{\;
\frac1{18}\sum_{w\in X_6^{\mathrm{cyc}}}\iota_C(w)\iota_C(w)^{\mathsf T}
=
\frac89\,I_3.\;}
\]
并且对任意 \(v\in\mathbb R^3\) 有
\[
\boxed{\;
\frac1{18}\sum_{w\in X_6^{\mathrm{cyc}}}\langle \iota_C(w),v\rangle^2
=
\frac89\,\|v\|_2^2.\;}
\]
等价地，对任意实对称矩阵 \(A\in M_3(\mathbb R)\)，若记 \(Q_A(v):=v^{\mathsf T}Av\)，则
\[
\boxed{\;
\sum_{w\in L}Q_A(\iota_C(w))
=
\sum_{w\in S}Q_A(\iota_C(w))
=
8\,\operatorname{tr}(A).\;}
\]
因此尽管两层分别具有 \(6\) 与 \(12\) 个向量、半径分别为 \(2\) 与 \(\sqrt2\)，它们对全部二次可观测量的总贡献完全相等；特别地，cyclic 层在该根字典下具有严格的二阶各向同性。
\end{theorem}
\begin{proof}
由表 \ref{tab:fold6_b3c3_root_dictionary_c3}，集合 \(L\) 正是
\[
\{\pm 2e_1,\ \pm 2e_2,\ \pm 2e_3\}.
\]
于是
\[
\sum_{w\in L}\iota_C(w)\iota_C(w)^{\mathsf T}
=
2\cdot(2e_1)(2e_1)^{\mathsf T}
+2\cdot(2e_2)(2e_2)^{\mathsf T}
+2\cdot(2e_3)(2e_3)^{\mathsf T}
=
8I_3.
\]

对集合 \(S\)，同一张表给出三组短根：
\[
\{(\pm1,\pm1,0)\},
\qquad
\{(\pm1,0,\pm1)\},
\qquad
\{(0,\pm1,\pm1)\}.
\]
对第一组有
\[
\sum_{(\varepsilon_1,\varepsilon_2)\in\{\pm1\}^2}
(\varepsilon_1 e_1+\varepsilon_2 e_2)
(\varepsilon_1 e_1+\varepsilon_2 e_2)^{\mathsf T}
=
4\bigl(e_1e_1^{\mathsf T}+e_2e_2^{\mathsf T}\bigr),
\]
因为所有混合项都按符号对称相消。其余两组同理，分别给出
\[
4\bigl(e_1e_1^{\mathsf T}+e_3e_3^{\mathsf T}\bigr),
\qquad
4\bigl(e_2e_2^{\mathsf T}+e_3e_3^{\mathsf T}\bigr).
\]
三式相加即得
\[
\sum_{w\in S}\iota_C(w)\iota_C(w)^{\mathsf T}=8I_3.
\]
总和公式随即成立。两边再除以 \(18\)，便得到归一化二次矩公式。

对任意 \(v\in\mathbb R^3\)，由
\[
\langle \iota_C(w),v\rangle^2
=
v^{\mathsf T}\bigl(\iota_C(w)\iota_C(w)^{\mathsf T}\bigr)v
\]
与总和公式
\[
\sum_{w\in X_6^{\mathrm{cyc}}}\iota_C(w)\iota_C(w)^{\mathsf T}=16I_3
\]
可得
\[
\sum_{w\in X_6^{\mathrm{cyc}}}\langle \iota_C(w),v\rangle^2
=
v^{\mathsf T}(16I_3)v
=
16\|v\|_2^2.
\]
再除以 \(18\) 即得
\[
\frac1{18}\sum_{w\in X_6^{\mathrm{cyc}}}\langle \iota_C(w),v\rangle^2
=
\frac89\,\|v\|_2^2.
\]

最后，对任意实对称矩阵 \(A\)，利用
\[
Q_A(v)=\operatorname{tr}\!\bigl(A\,vv^{\mathsf T}\bigr)
\]
与迹的线性性，可得
\[
\sum_{w\in L}Q_A(\iota_C(w))
=
\operatorname{tr}\!\left(
A\sum_{w\in L}\iota_C(w)\iota_C(w)^{\mathsf T}
\right)
=
\operatorname{tr}(A\cdot 8I_3)
=
8\,\operatorname{tr}(A),
\]
而 \(S\) 的情形完全相同。证毕。
\end{proof}

\begin{corollary}[window-\texorpdfstring{$6$}{6} 的 \texorpdfstring{$C_3$}{C3} 根字典不存在非零一次不变量]\label{cor:xi-window6-c3-no-nonzero-linear-invariant}
沿用定理 \ref{thm:xi-window6-c3-quadratic-energy-equipartition} 的记号，则
\[
\sum_{w\in X_6^{\mathrm{cyc}}}\iota_C(w)=0.
\]
因此对任意线性泛函 \(\ell:\mathbb R^3\to\mathbb R\)，都有
\[
\sum_{w\in X_6^{\mathrm{cyc}}}\ell\bigl(\iota_C(w)\bigr)=0.
\]
\end{corollary}
\begin{proof}
由表 \ref{tab:fold6_b3c3_root_dictionary_c3}，\(X_6^{\mathrm{cyc}}\) 中的每个向量都与其相反向量成对出现：长根部分为
\[
\{\pm(2,0,0),\ \pm(0,2,0),\ \pm(0,0,2)\},
\]
短根部分为
\[
\{(\pm1,\pm1,0),\ (\pm1,0,\pm1),\ (0,\pm1,\pm1)\}.
\]
故逐对相加即得
\[
\sum_{w\in X_6^{\mathrm{cyc}}}\iota_C(w)=0.
\]
对任意线性泛函 \(\ell\)，再由线性性立得
\[
\sum_{w\in X_6^{\mathrm{cyc}}}\ell\bigl(\iota_C(w)\bigr)
=
\ell\!\left(\sum_{w\in X_6^{\mathrm{cyc}}}\iota_C(w)\right)
=
0.
\]
证毕。
\end{proof}
\leanverified{paper\_xi\_window6\_c3\_no\_nonzero\_linear\_invariant}

\endinput
